% ====================================================================
%  CHAPTER: INTRODUCTION
% ====================================================================
\chapter{Introduction}
\label{ch:introduction}

% ====================================================================
%  §1  MOTIVATIONS
% ====================================================================
\section{Motivations}
\label{sec:intro-motivations}
Deep neural networks (DNNs) have become the dominant computational
paradigm across a wide range of applications, from image classification
and object detection to natural language processing and autonomous
driving~\cite{lecun2015deep}.
The computational cost of modern DNNs, however, is immense: a single
inference pass of a large convolutional network such as VGG16 requires
billions of multiply-and-accumulate (MAC) operations, while training
demands orders of magnitude more~\cite{vgg16}.
While general-purpose processors and graphics processing units (GPUs)
deliver high computational throughput, their rigid memory hierarchies
often struggle to meet the stringent energy-efficiency (performance-per-watt)
requirements of specialized edge devices. This bottleneck has motivated the development of dedicated
\emph{spatial architectures} (SAs): custom accelerators that exploit
massive data-level parallelism through arrays of processing elements (PEs)
interconnected by customized on-chip memory networks.


The classic execution model for a neural network on a SA is the layer-by-layer execution model, where the neural network is computed strictly one layer at a time. This approach introduces a severe constraint: the output activation tensor must be entirely written to off-chip DRAM after one layer completes, only to be read back when the next layer begins. This creates a
``memory wall'': off-chip DRAM accesses are roughly two orders of magnitude more expensive than on-chip SRAM accesses in terms of energy. This traffic pattern typically dominates the total energy budget for activation-heavy
workloads. Furthermore, the limited bandwidth of DRAM frequently leads to
processing stalls when handling heavy read and write traffic of intermediate data, severely
degrading latency.

\emph{Layer fusion} has emerged as a promising strategy to address this
bottleneck by computing multiple consecutive layers in a single, interleaved
schedule. By tiling the execution, intermediate activations
can be produced and consumed entirely within on-chip buffers in localized
chunks, effectively eliminating the round trip through external DRAM.
Hardware implementations such as DepFiN~\cite{goetschalckx2023depfin}
have demonstrated that depth-first fused execution can achieve
significant energy savings and latency reductions for certain workloads.

However, a crucial challenge remains even for single-layer execution: extracting optimal performance and energy efficiency from these hardware designs relies critically on the \emph{mapping}, a schedule that determines the order in which loop iterations are executed, how data is tiled and distributed across PEs, and what data is retained at each level of the memory hierarchy. Even for a single convolutional layer on a moderately complex SA, the set of all possible mappings (the map space) exceeds $10^{19}$ candidates. This staggering complexity has driven the development of fast analytical cost models and automated mappers to efficiently search for optimal single-layer schedules.
Furthermore, fusing $L$~layers exponentially multiplies the number of loop dimensions, introduces strict producer--consumer dependencies, and creates new opportunities for data reuse. This dimensional proliferation causes an astronomical growth in the map space.  For instance, an 10-layer fusion of FSRCNN on the DepFiN architecture expands the map space to over $10^{515}$ candidates, which is more than 400 orders of magnitude larger than its single-layer equivalent~\cite{Factorflow}. To tame this intractable space, it is fundamental to study the cascading tile-size dependencies between producer and consumer layers. Enforcing these dependencies can rigorously prune the map space of physically infeasible mappings that an automated mapper would otherwise needlessly evaluate during the search for optimality. Moreover, maximizing data reuse requires optimizing not only the intra-layer (within a single layer) schedules but also the inter-layer tile mappings, necessitating a systematic analysis of the different types of buffered and immediate reuse mechanisms that fused execution unlocks.

Despite the increase in complexity, since the mapping of a single layer is natively expressed as a loop nest, it is highly desirable to maintain this representational consistency across the extended map space. Describing the fused-layer mapping as a single, unified loop nest via the Extended Einsum notation provides an exact, coherent mathematical representation. While existing fused-layer exploration frameworks, such as DeFiNES~\cite{mei2023defines} and LoopTree~\cite{gilbert2024looptree}, have made significant progress in exploring the fused dataflow design space, no prior tool represents both single-layer and fused-layer mappings under a unified, coherent loop-nest abstraction.

Therefore, the core motivation of this thesis is to formally explore the fused-layer dataflow design space through a unified analytical cost model. This work seeks to rigorously study the architectural pros and cons of implementing fused layers, analyzing how different hardware memory hierarchies and workload characteristics (such as activation versus weight dominance) dictate its success or failure.

% ====================================================================
%  §2  OBJECTIVES
% ====================================================================
\section{Objectives}
\label{sec:intro-objectives}

The primary objective of this study is to formally explore the design space
of multi-layer layer fusion and quantify its inherent performance and energy
trade-offs. To navigate this space, our thesis extends
the FactorFlow analytical mapping framework for single layer. The specific objectives are to:
%
\begin{enumerate}
  \item \textbf{Formalize the Fused-Layer Design Space:} Characterize the
        theoretical complexities of layer fusion, including dimension proliferation,
        operand residency rules, and the cascading inter-layer tile dependencies
        that dictate the reuse-recompute trade-off.

  \item \textbf{Develop Constraint-Based Pruning:} Tame the astronomical
        map-space explosion by establishing an exact, scheduling pruning
        strategy that enforces strict producer--consumer tile-size inequalities at
        every level of the hardware hierarchy.

  \item \textbf{Construct an Analytical Evaluation Engine:} Extend FactorFlow's
        existing factor-based abstraction and extended Einsum notation to unify
        multi-layer loop nests (via aliasing and decoupling) without sacrificing
        intra-layer mapping independence. Generalize the single-layer cost model to assess
        memory operations, energy, and latency on a per-layer basis.

  \item \textbf{Validate the Model:} Validate the extended analytical framework
        against a published hardware baseline (DepFiN), demonstrating that the
        predictions for complex, deep fusion topologies match the reference.

  \item \textbf{Quantify the Pros and Cons of Fusion:} Conduct a comprehensive
        design-space exploration across different spatial architectures (DepFiN-Like
        and Eyeriss-Like) and CNN workloads (activation-dominant vs. weight-dominant).
        By comparing full fusion, partial fusion, and single-layer execution modes,
        this objective aims to definitively isolate the architectural and algorithmic
        conditions under which layer fusion provides a tangible benefit versus
        a detrimental overhead.
\end{enumerate}

% ====================================================================
%  §3  THESIS ORGANIZATION
% ====================================================================
\section{Thesis Organization}
\label{sec:intro-organisation}

The remainder of this thesis is organized as follows:

\begin{description}
  \item[Chapter~\ref{ch:background} -- Background]
    introduces the foundational concepts: artificial intelligence
    and deep learning, DNN workloads and the extended Einsum notation,
    hardware acceleration with spatial architectures, the mapping
    problem and map-space explosion, analytical cost modeling,
    and the mechanics of layer fusion.

  \item[Chapter~\ref{ch:fused-design-space} -- Fused-Layer Mapping Design Space]
    formalizes the core theoretical challenges of the fused-layer mapping problem:
    the exponential proliferation of dimensions, cascading tile dependencies,
    operand categorization, memory bypass rules, and the fundamental
    reuse-recompute trade-offs that dictate map-space feasibility.

   \item[Chapter~\ref{ch:Related Works} -- Related Works]
    examines the state of the art in layer-by-layer dataflow
    exploration tools (Timeloop, FactorFlow) and fused-layer 
    evaluation frameworks (DeFiNES, LoopTree), contextualizing 
    them within the broader landscape of layer fusion methodologies.


  \item[Chapter~\ref{ch:proposed-solution} -- Proposed Solution]
    presents the analytical framework developed to explore the design space:
    the aliasing and decoupling abstraction, the exact constraint-based
    pruning, the generalized cost model, and
    software implementation details.

  \item[Chapter~\ref{ch:validation} -- Validation]
    validates the extended analytical model against the DepFiN accelerator,
    analyzing the agreement between the model's predicted data-movement and energy metrics and the published figures from a baseline fused workload experiment.

  \item[Chapter~\ref{ch:experimental-results} -- Experimental Results]
    presents the comprehensive design-space exploration. It details the setup,
    single-architecture sensitivity analyses and evaluates the critical pros and cons of full, partial, and single-layer execution across varied architectural constraints.

  \item[Chapter~\ref{ch:conclusions} -- Conclusions and Future Work]
    summarizes the theoretical and empirical findings of the study and
    outlines directions for future research toward automated heuristic mappers.
\end{description}